{\footnotesize
\begin{tabularx}{\textwidth}{>{\fontSemiBold\color{BrandNavy}\hsize=0.503\hsize\sloppy\arraybackslash}X C{1.70cm} >{\hsize=1.497\hsize\sloppy\arraybackslash}X}
\arrayrulecolor{GithubTableBorder}\toprule
\rowcolor{BrandTableHeader}
\textbf{Pergunta de avaliação} & \textbf{Seção Revisada} & \textbf{Nova Situação e Evidência Verificável (Versão 2.0)} \\
\midrule
\arrayrulecolor{BrandGrayRule}
\textbackslash1. O escopo e as entregas do projeto estão claros? & \texttt{4.1} & \textbf{Descrito.} A seção 4.1 agora inclui o diagrama da EAP em 3 níveis e a tabela das \textbf{3 Macroentregas} (\texttt{v1.0.0}, \texttt{v2.0.0} e \texttt{v3.0.0}), detalhando exatamente os pacotes entregues em produção e a evolução comercial. \\
\midrule
\textbackslash2. As atividades possuem responsáveis e prazos? & \texttt{4.1} & \textbf{Descrito.} Tanto a tabela de Marcos Comerciais quanto a tabela das 3 Sprints de execução especificam nominalmente os engenheiros responsáveis (\texttt{Gabriel\ Lourenço}, \texttt{Carlos\ Eduardo}, \texttt{Kayque\ Vicente}, \texttt{Yury\ Nogueira}) e as datas de início e término. \\
\midrule
\textbackslash3. Os riscos e suas respostas estão definidos? & \texttt{4.2} & \textbf{Descrito.} A seção 4.2 explicita os critérios de Probabilidade, Impacto e Prioridade, o ciclo de reavaliação por Sprint e a matriz com 7 riscos reais de engenharia/negócio, respostas técnicas e responsáveis nominais. \\
\midrule
\textbackslash4. Está definido como acompanhar o andamento e tratar atrasos? & \texttt{4.1} e \texttt{4.5} & \textbf{Descrito.} A seção 4.1 formaliza a subseção \emph{``Rotina de Acompanhamento e Tratamento de Desvios e Atrasos''} (gatilho de 48h e pareamento técnico), complementada na seção 4.5 por métricas com metas objetivas e resultados aferidos (\texttt{Lead\ Time\ \textless{}=\ 48h}, \texttt{WIP\ \textless{}=\ 2}, \texttt{770\ testes}). \\
\arrayrulecolor{GithubTableBorder}\bottomrule
\end{tabularx}%
}
